% =======================================================================
% Playing It Safe: Managerial Speech Uncertainty and
%                  Corporate Financial Conservatism
% -----------------------------------------------------------------------
% Draft skeleton, immature first version.
% Template: docs/Draft/DraftTemplate.txt (strictly followed).
% Decisions log: docs/Draft/DECISIONS.md
% Authoritative tables: outputs/all_tables.tex
% Authoritative findings: outputs/findings.txt
% =======================================================================

\documentclass[11pt,a4paper,twocolumn]{article}

% ---- Geometry and typography ----
% Compile with xelatex (fontspec requires a Unicode engine for system
% TTF/OTF fonts such as the Windows Times New Roman).
\usepackage[margin=0.9in,columnsep=0.32in]{geometry}
\usepackage[no-math]{fontspec}
\setmainfont{Times New Roman}
\usepackage{newtxmath}   % Times-compatible math companion
\usepackage{setspace}
\setlength{\parindent}{1.5em}
\setlength{\parskip}{0.25em}

% ---- Math and tables ----
\usepackage{amsmath}
% NOTE: amssymb intentionally omitted --- newtxmath already provides
% the blackboard-bold symbols and loading amssymb alongside newtxmath
% causes a \Bbbk redefinition clash.
\usepackage{booktabs}
\usepackage{graphicx}
\usepackage{float}

% ---- Citations, links, quoting ----
\usepackage{natbib}
\usepackage[colorlinks=true,citecolor=blue,linkcolor=blue,urlcolor=blue]{hyperref}
\usepackage{csquotes}

% ---- Draft-stage placeholder markers ----
\usepackage{xcolor}
\definecolor{placeholderaccent}{RGB}{204,85,0}
\newenvironment{placeholder}{%
  \par\medskip\noindent%
  {\color{placeholderaccent}\rule{\linewidth}{0.4pt}}\par%
  \noindent\textbf{\color{placeholderaccent}[PLACEHOLDER]}\par%
  \smallskip\itshape%
}{%
  \par\upshape\smallskip%
  {\color{placeholderaccent}\rule{\linewidth}{0.4pt}}\par\medskip%
}

% ---- Title block ----
\title{%
  Playing It Safe:\\[0.3em]
  \Large Managerial Speech Uncertainty and\\ Corporate Financial Conservatism
}

\author{%
  [Author Name]\\
  [Institution]\\
  \texttt{[email@example.com]}
  \and
  [Author Name]\\
  [Institution]\\
  \texttt{[email@example.com]}
  \and
  [Author Name]\\
  [Institution]\\
  \texttt{[email@example.com]}
}

\date{[Month Year]}

% =======================================================================
\begin{document}
% =======================================================================

% ---- Title + abstract + keywords/JEL span both columns (full width) ----
\makeatletter
\twocolumn[%
  \begin{@twocolumnfalse}
    \maketitle
    \begin{abstract}
      \noindent
      \begin{placeholder}
      ABSTRACT (target 150--200 words). Must cover:
      (i) Measurement of managerial speech uncertainty from earnings
      conference calls using the DWZ (2021) persistent-style framework,
      decomposed into answer-side (Q\&A, \texttt{UncAnsMgr}) and
      presentation-side (scripted, \texttt{UncPreMgr}) components, and
      extended beyond the CEO to the full management team.
      (ii) Main finding: speech uncertainty predicts corporate financial
      conservatism across cash holdings, leverage, dividends, and buybacks.
      (iii) Defensive-investment result: capex \emph{rises} under manager
      uncertainty and is \emph{amplified} by product-market competition ---
      a competitive real-options story rather than classical wait-and-see
      delay.
      (iv) Sample: 112{,}968 earnings calls, 2{,}429 firms, 2002--2018
      (Compustat, CapIQ, Execucomp, CRSP).
      (v) External validation: the measure responds to macro economic
      policy uncertainty, political risk, and product-market competition
      --- it is not noise.
      \end{placeholder}
    \end{abstract}

    \medskip
    \noindent\textbf{Keywords:}
    \begin{placeholder}
    Keywords --- managerial speech uncertainty; earnings
    conference calls; precautionary cash holdings; financial conservatism;
    competitive real options; DWZ 2021.
    \end{placeholder}

    \medskip
    \noindent\textbf{JEL Classification:}
    \begin{placeholder}
    JEL codes --- G30 (General Corporate Finance); G32
    (Financing Policy; Capital Structure); G35 (Payout Policy); D81
    (Criteria for Decision-Making under Risk and Uncertainty); M41
    (Accounting).
    \end{placeholder}
    \vspace{1.5em}
  \end{@twocolumnfalse}
]
\makeatother

% =======================================================================
% Section I -- Introduction
% =======================================================================
\section{Introduction}
\label{sec:intro}

\begin{placeholder}
INTRODUCTION must cover, in order:
\begin{enumerate}
  \item \textbf{Hook.} Why managerial communication on earnings calls
        matters for corporate policy: investors, analysts, and the
        disclosure literature all treat the call as a signal, but the
        specific content of \emph{uncertainty} in managerial language
        has been measured almost exclusively for the CEO.
  \item \textbf{Gap.} Prior work (DWZ 2021 and followers) focuses on CEO
        and CFO language separately and on the full call; we extend the
        measurement to the broader management team (CFO + non-CEO
        speakers in the Q\&A) and decompose the measure into
        answer-side (Q\&A, spontaneous) and presentation-side
        (scripted) components.
  \item \textbf{Central claim.} Managerial speech uncertainty predicts
        corporate financial conservatism: cash holdings rise, leverage
        drifts down at the one-quarter lead, dividends fall, share
        repurchases fall.
  \item \textbf{Narrative tension / secondary claim.} Capital expenditure
        \emph{rises} with manager uncertainty and the effect is
        \emph{amplified} by product-market competition. We frame this
        as a competitive real-options result in the spirit of
        Grenadier (2002) and Aguerrevere (2009): when waiting risks
        preemption, uncertain managers invest \emph{defensively} to
        preserve market position.
  \item \textbf{Data footprint.} 112{,}968 earnings conference calls,
        2{,}429 unique firms, 2002 Q1--2018 Q4. Sources: CapIQ call
        transcripts and speaker metadata, Compustat quarterly
        fundamentals, CRSP daily prices, Execucomp CEO/CFO flags.
  \item \textbf{External validation.} The measure is not noise: it
        responds to macro economic policy uncertainty (Baker-Bloom-Davis),
        political risk (Hassan et al.), and product-market competition
        (Hoberg-Phillips TSIMM).
  \item \textbf{Contribution.} (a) Decomposition of managerial
        speech uncertainty into answer vs.\ presentation channels;
        (b) extension beyond CEO to full management; (c) identification
        of multiple conservatism channels inside a single firm-quarter
        panel; (d) the defensive-investment result as a contribution to
        the uncertainty--investment literature.
  \item \textbf{Roadmap.} Section~II lays out the framework, measurement,
        and estimation. Section~III presents the main results.
        Section~IV collects additional analyses, validity tests, and
        null boundary results. Section~V concludes.
\end{enumerate}
\end{placeholder}

\clearpage

% =======================================================================
% Section II -- Conceptual Framework and Empirical Strategy
% =======================================================================
\section{Conceptual Framework and Empirical Strategy}
\label{sec:framework}

\subsection{Conceptual Framework}
\label{sec:framework_concept}
\begin{placeholder}
Frame \emph{financial conservatism} as the core dependent construct.
Draw on four literatures:
(i) precautionary cash holdings (Bates, Kahle, \& Stulz 2009;
Opler, Pinkowitz, Stulz, \& Williamson 1999);
(ii) capital structure and leverage adjustment (Myers 1984 pecking order;
Frank \& Goyal 2003);
(iii) payout smoothing and dividend commitment
(Lintner 1956; Brav, Graham, Harvey, \& Michaely 2005);
(iv) real options and uncertainty--investment
(Dixit \& Pindyck 1994; Bloom 2014 --- note: asymmetric payoffs,
\emph{not} ``explore/resolve uncertainty'' per memory correction).
The novelty is the IV, not the DVs: we use a new measure of manager
speech uncertainty to identify the conservatism response in a single
firm-quarter panel.
\end{placeholder}

\subsection{Theory: Managerial Speech Uncertainty as a Persistent Style}
\label{sec:framework_theory}
\begin{placeholder}
Build on \citet{DWZ2021} persistent managerial style. Key points:
\begin{itemize}
  \item Managers exhibit persistent linguistic styles across calls.
  \item We decompose uncertainty into answer-side (\texttt{UncAns})
        and presentation-side (\texttt{UncPre}) components, mirroring
        the scripted/unscripted distinction in the calls.
  \item We extend from CEO-only to a full management grouping
        (\texttt{UncAnsMgr}, \texttt{UncPreMgr}) that pools the CFO and
        other non-CEO executives speaking on the call.
  \item Procedural precedent for \emph{management}-level measurement:
        Bushee, Gow, \& Taylor (2018) --- via Ian Gow's
        \texttt{iangow/bgt} replication code.
  \item Procedural precedent for CFO measurement: Larcker \& Zakolyukina
        (2012), using the FactSet XML title field.
\end{itemize}
\end{placeholder}

\subsection{Estimation of Main Variables}
\label{sec:framework_estimation}
\begin{placeholder}
Describe the linguistic uncertainty dictionary, normalization, and the
manager-vs-CEO speaker classification pipeline:
\begin{itemize}
  \item Loughran--McDonald uncertainty word list, call-share
        normalization, DWZ 2021 style-residual persistence component.
  \item Speaker classification: Execucomp \texttt{ceoann}/\texttt{cfoann}
        annual flags merged on gvkey--year; manager classifier picks up
        remaining executives from transcript metadata.
  \item \textbf{Known blocker:} current \texttt{UncAnsMgr} construction
        is broken in the pipeline (memory
        \texttt{project\_step22\_schema\_bug}); CFO / ManagerV2 rebuild
        plan is finalized (memory \texttt{project\_cfo\_mgrv2\_plan\_v2})
        and ready to execute. Results in the tables reflect the
        \emph{current} working build; we will rerun after the rebuild.
\end{itemize}
\end{placeholder}

\subsection{Methodology and Empirical Design}
\label{sec:framework_method}
\begin{placeholder}
\begin{itemize}
  \item \textbf{Panel.} Firm-quarter, 2002 Q1--2018 Q4. Main sample
        excludes financial firms (SIC 6000--6999) and utilities
        (SIC 4900--4999).
  \item \textbf{Estimator.} PanelOLS via Python \texttt{linearmodels}.
        Two-way clustering at the firm and calendar-year-quarter level.
        Calendar year-quarter fixed effects (not fiscal) per
        memory \texttt{feedback\_calendar\_yr\_qtr\_fe}.
  \item \textbf{12-column specification ladder.} Three specifications
        (Industry+Year, Industry+Year+Extended Controls, Firm+Year,
        Firm+Year+Extended Controls, Year-Quarter FE variants) applied
        contemporaneously and at the one-quarter lead, giving 12 columns
        per main-table DV.
  \item \textbf{Lagged DV as base control} in every column (memory
        \texttt{feedback\_lagged\_dv}).
  \item \textbf{Inference.} IVs tested one-tailed when direction is
        predicted (financial-conservatism hypotheses); controls
        two-tailed.
\end{itemize}
\end{placeholder}

\subsection{Specification and Measurement of Key Constructs}
\label{sec:framework_specs}
\begin{placeholder}
\textbf{Base controls:} \texttt{lnAssets}, \texttt{TobinsQ},
\texttt{ROA}, \texttt{Leverage}, \texttt{Capex}, \texttt{CashRatio},
\texttt{DivDummy}, \texttt{sCFO}, \texttt{Lagged\_DV}.\\
\textbf{Extended controls:} \texttt{SalesGrowth}, \texttt{RDSales},
\texttt{CashFlowAt}, \texttt{DailyVola}.\\
\textbf{Bad-control exclusions:} per DV, we drop controls that share
the dependent variable's numerator concept (memory
\texttt{feedback\_bad\_control\_exclusion}). For example, when
\texttt{CashRatio} is the DV we do not use \texttt{CashFlowAt} and
\texttt{CashRatio} itself as right-hand-side controls; the full map is
enforced by the shared runner configuration.
\end{placeholder}

\clearpage

% =======================================================================
% Section III -- Main Empirical Analyses
% =======================================================================
\section{Main Empirical Analyses}
\label{sec:main}

\subsection{Data, Sample, and Variable Construction}
\label{sec:data}
\begin{placeholder}
\textbf{Sample.} 112{,}968 earnings conference calls, 2{,}429 unique
firms, 2002 Q1--2018 Q4. Build pipeline:
\begin{itemize}
  \item CapIQ call transcripts and structured speaker metadata;
        CFO/non-CEO manager flags via Execucomp \texttt{ceoann}/
        \texttt{cfoann} annual fields merged on gvkey--year.
  \item Compustat quarterly fundamentals for financial-statement items.
  \item CRSP daily prices and volume for microstructure tests in
        Appendix~II.
  \item Main sample excludes financial (SIC 6000--6999) and utility
        (SIC 4900--4999) firms.
\end{itemize}
\textbf{Table~1 Summary Statistics} placeholder --- generate from
\texttt{findings.txt} or a dedicated summary script.
\end{placeholder}

% -----------------------------------------------------------------------
\subsection{Precautionary Liquidity: Cash Holdings and the Rating Channel}
\label{sec:liquidity}
\begin{placeholder}
\textbf{Subnarrative.} Uncertain managers hoard cash; financial
constraint --- proxied by being unrated --- amplifies the effect.
\begin{itemize}
  \item \textbf{H1 (Table~\ref{tab:h1}).} \texttt{UncAnsMgr}
        $+0.0034$** to $+0.0072$*** across the six contemporaneous
        specifications of \texttt{CashRatio}; lead horizon weaker.
  \item \textbf{H1.2 (Table~\ref{tab:h1_2}).} Three-category rating
        moderation: main effect \texttt{Manager\_QA\_Unc\_c}
        $+0.0055$*** to $+0.0057$***; unrated interaction
        \texttt{MgrQAUnc\_x\_Unrated} $+0.0068$*** to $+0.0073$***.
        Below-investment-grade interaction null.
  \item \textbf{Interpretation.} Classical precautionary savings motive
        \citep{BatesKahleStulz2009} amplified when external liquidity
        access is absent. The unrated firm is the marginal cash hoarder.
\end{itemize}
References: Tables~\ref{tab:h1}, \ref{tab:h1_2}.
\end{placeholder}

% -----------------------------------------------------------------------
\subsection{Capital-Structure Discipline: Book Leverage and Debt-to-Capital}
\label{sec:leverage}
\begin{placeholder}
\textbf{Subnarrative.} Deleveraging is sluggish. Nothing appears
contemporaneously, but a modest negative pull shows up at the
one-quarter lead.
\begin{itemize}
  \item \textbf{H4a (Table~\ref{tab:h4a}).} \texttt{UncAnsMgr} null
        contemporaneously; $-0.0050$* to $-0.0090$** at the lead,
        surviving firm and year-quarter fixed effects.
  \item \textbf{H4b (Table~\ref{tab:h4b}).} Mostly null; a single
        $-0.0154$*** at lead in one specification. Weakest of the two.
  \item \textbf{Honest caveat.} This is the weakest section of the
        main body. Frame as \emph{suggestive} deleveraging, not a
        headline. Discuss the pecking-order implication and flag that
        formal financing-choice margins (H19b, H20b in Appendix~I)
        are similarly fragile.
\end{itemize}
References: Tables~\ref{tab:h4a}, \ref{tab:h4b}.
\end{placeholder}

% -----------------------------------------------------------------------
\subsection{Distribution Restraint: Dividends, Payer Status, and Repurchases}
\label{sec:payout}
\begin{placeholder}
\textbf{Subnarrative.} When managers sound uncertain, cash stays in
the firm. Dividends fall, the probability of being a dividend payer
falls, and share repurchases fall. Combined, these three channels
show that \emph{total} distributions to shareholders contract.
\begin{itemize}
  \item \textbf{H12 (Table~\ref{tab:h12}).} Quarterly Payout Ratio.
        \texttt{UncPreMgr} $-0.0299$** to $-0.0529$*** ---
        presentation-side channel.
  \item \textbf{H12b (Table~\ref{tab:h12b}).} Dividend Payer Indicator
        (Hoberg-Prabhala 2009 analog, linear probability + firm FE).
        \texttt{UncPreMgr} $-0.0055$** to $-0.0113$***.
  \item \textbf{H17 (Table~\ref{tab:h17}).} Repurchase Intensity
        (de-cumulated \texttt{prstkcy}). \texttt{UncAnsMgr} $-0.0005$* to
        $-0.0008$*** contemporaneously, with \texttt{UncPreMgr} showing
        a small \emph{positive} partial correlation ($+0.0005$* to
        $+0.0007$**) --- a partial cross-channel tension flagged below.
  \item \textbf{Channel note.} Payouts run through presentation-side
        uncertainty (\texttt{UncPreMgr}); repurchases run through
        answer-side uncertainty (\texttt{UncAnsMgr}). Whether to
        feature this as a two-channel contribution or smooth it over
        is an open decision (see DECISIONS.md \S3.3).
\end{itemize}
References: Tables~\ref{tab:h12}, \ref{tab:h12b}, \ref{tab:h17}.\\
\textbf{Open:} H17 placement pending user confirmation --- current
proposal groups it here with payouts under total distribution.
\end{placeholder}

% -----------------------------------------------------------------------
\subsection{Defensive Investment: Capital Expenditure and the Competition Channel}
\label{sec:investment}
\begin{placeholder}
\textbf{Subnarrative.} Capex \emph{rises} under manager uncertainty,
and the effect is \emph{amplified} by product-market competition. This
is \textbf{the central narrative tension} with the ``playing it safe''
title: on the surface, more investment looks like more risk-taking.
We reconcile it as a competitive real-options result: when waiting
risks preemption, the option value of delay collapses, and the
uncertain manager invests defensively to preserve competitive position.
\begin{itemize}
  \item \textbf{H13 (Table~\ref{tab:h13}).} \texttt{UncAnsMgr}
        $+0.0020$*** to $+0.0049$*** in the extended-control
        specifications.
  \item \textbf{H13.1 (Table~\ref{tab:h13_1}).} Product-similarity
        moderation: main effect $+0.0010$** to $+0.0017$**, interaction
        \texttt{MgrQAUnc\_x\_zlogTSIMM} $+0.0011$* to $+0.0024$**
        (positive --- competition amplifies the effect).
  \item \textbf{Theoretical framing.} Classical real options
        \citep{DixitPindyck1994} predicts \emph{less} investment under
        uncertainty. We instead see more. The competitive real-options
        literature \citep{Grenadier2002,Aguerrevere2009} rationalizes
        this: in a setting where rivals may preempt, waiting has
        negative strategic value, and the optimal response to higher
        uncertainty can be \emph{earlier} exercise. Our H13.1 interaction
        sign is the empirical signature of this mechanism.
  \item \textbf{Persistence.} The effect is not a one-quarter blip. See
        Appendix~\ref{sec:appI} (H13.2) for lead horizons $t+1$ through
        $t+4$.
\end{itemize}
References: Tables~\ref{tab:h13}, \ref{tab:h13_1}.\\
\textbf{Open:} H13 framing is the central open decision (DECISIONS.md
\S3.1). Options: (a) defensive-hedge framing (current proposal);
(b) split thesis, two independent stories; (c) demote H13 out of
main.
\end{placeholder}

% -----------------------------------------------------------------------
\subsection{Non-Distinguishing Investment: Research and Development}
\label{sec:rd}
\begin{placeholder}
\textbf{Subnarrative.} R\&D does not respond to managerial speech
uncertainty in either direction.
\begin{itemize}
  \item \textbf{H16 (Table~\ref{tab:h16}).} \texttt{RDSales} (Jiang 2021
        construction, \texttt{xrdy}/\texttt{saley}). All four IVs are
        statistically null across all 12 specifications.
  \item \textbf{Interpretation.} R\&D is a long-horizon, partly sticky
        commitment: prior-period research projects, accounting treatment
        under ASC 730, and staffing contracts all buffer R\&D from
        short-run communication signals. Absence of response is
        itself informative --- it distinguishes R\&D from capex
        (H13) and suggests the defensive-investment mechanism is
        tangible-capacity specific.
\end{itemize}
Reference: Table~\ref{tab:h16}.\\
\textbf{Open:} H16 placement pending user decision (DECISIONS.md
\S3.2). Options: keep here as ``no response'' or relocate to
Appendix~I.
\end{placeholder}

\clearpage

% =======================================================================
% Section IV -- Additional Analyses
% =======================================================================
\section{Additional Analyses}
\label{sec:additional}

% =======================================================================
% Section IV.A -- Appendix I: Robustness, Channels, and Horizons
% =======================================================================
\subsection{Robustness, Channels, and Horizons (Appendix I)}
\label{sec:appI}

\subsubsection{Product-Similarity Channel for Cash (Null)}
\label{sec:appI_h11}
\begin{placeholder}
\textbf{Subnarrative.} Competition does \emph{not} moderate the cash
story --- the channel that works for capex (H13.1) does not work for
precautionary liquidity.
\begin{itemize}
  \item \textbf{H1.1 (Table~\ref{tab:h1_1}).} Main
        \texttt{Manager\_QA\_Unc\_c} $+0.0045$***;
        \texttt{z\_log\_TotalSimilarity} $+0.0057$***; interaction
        \texttt{MgrQAUnc\_x\_zlogTSIMM} $+0.0017$ (null).
  \item \textbf{Binary variant H1.1b (Table~\ref{tab:h1_1b}).}
        \texttt{MgrQAUnc\_x\_HighTSIMM} $+0.0010$ (null).
  \item \textbf{Interpretation.} Precautionary cash hoarding is a
        general response to managerial uncertainty, not a
        competition-specific one. The \emph{investment} decision is
        where competitive dynamics bite (H13.1).
\end{itemize}
References: Tables~\ref{tab:h1_1}, \ref{tab:h1_1b}.
\end{placeholder}

\subsubsection{Capex Persistence: Lead Horizons}
\label{sec:appI_h132}
\begin{placeholder}
\textbf{Subnarrative.} The defensive-investment effect is \emph{not}
a one-quarter blip. Capex responds at $t+1$ and the response
\emph{grows} through $t+4$.
\begin{itemize}
  \item \textbf{H13.2 (Table~\ref{tab:h13_2}).} \texttt{UncAnsMgr}
        on \texttt{Capex} at:\\
        $t+1$: $+0.0049$*** (Ind.\ FE spec), $+0.0028$*** (Firm FE spec).\\
        $t+2$: $+0.0036$*** / $+0.0035$***.\\
        $t+3$: $+0.0070$*** / $+0.0049$***.\\
        $t+4$: $+0.0071$*** / $+0.0048$***.
  \item \textbf{Interpretation.} The competitive-preemption
        interpretation requires a sustained commitment to capacity,
        not a single project. The horizon results are consistent with
        that: uncertain managers commit to multi-quarter capex
        programs.
\end{itemize}
Reference: Table~\ref{tab:h13_2}.
\end{placeholder}

\subsubsection{Financing-Mix Margins (Weak)}
\label{sec:appI_h19b}
\begin{placeholder}
\textbf{Subnarrative.} Direct financing-choice margins point in the
conservative direction but are statistically fragile.
\begin{itemize}
  \item \textbf{H19b (Table~\ref{tab:h19b}).} External vs.\ internal
        financing, Chang, Dasgupta, \& Hilary (2006) cash-flow
        debt-classification. \texttt{UncAnsMgr} $-0.0155$* to
        $-0.0160$* at the lead --- suggestive shift toward internal
        financing, consistent with the lagged deleveraging in H4a/b.
  \item \textbf{H20b (Table~\ref{tab:h20b}).} Debt vs.\ equity choice,
        same CDH (2006) framework. Mostly null; \texttt{UncPreCEO}
        $-0.0269$* to $-0.0320$** in two specifications.
  \item \textbf{Honest framing.} Directions are consistent with
        conservatism, but statistical strength is weak. These results
        reinforce rather than establish the capital-structure
        section in the main body.
\end{itemize}
References: Tables~\ref{tab:h19b}, \ref{tab:h20b}.
\end{placeholder}

% =======================================================================
% Section IV.B -- Appendix II: Validity, Drivers, and Boundary Tests
% =======================================================================
\subsection{Validity, Drivers, and Boundary Tests (Appendix II)}
\label{sec:appII}

\subsubsection{External Drivers of Managerial Speech Uncertainty}
\label{sec:appII_drivers}
\begin{placeholder}
\textbf{Subnarrative.} The measure is not noise. Managerial speech
uncertainty responds to three independent external forcings ---
political risk, economic policy uncertainty, and product-market
competition --- all in the expected direction. This subsection
validates the IV.
\begin{itemize}
  \item \textbf{H11 (Table~\ref{tab:h11}).} Contemporaneous
        Hassan et al.\ PRisk $\to$ uncertainty language:
        \texttt{UncPreMgr} coefficient $+0.1174$***;
        \texttt{UncPreCEO} coefficient $+0.1078$***.
  \item \textbf{H11-Lag (Table~\ref{tab:h11_lag}).} Lagged PRisk
        variant: $+0.1219$*** / $+0.1122$***, robust to both FE
        schemes.
  \item \textbf{H23 (Table~\ref{tab:h23}).} Reverse-direction test of
        Hoberg-Phillips TSIMM $\to$ managerial uncertainty:
        \texttt{UncPreMgr} $+0.1943$***; \texttt{UncPreCEO}
        $+0.2491$***. Product-market competition increases hedged
        language.
  \item \textbf{H24 (Table~\ref{tab:h24}).} $\log(\text{US EPU})_t \to$
        uncertainty: coefficient range $+0.0117$* to $+0.0239$*** across
        all eight specifications.
  \item \textbf{H24b (Table~\ref{tab:h24b}).} $\log(\text{GEPU})_t
        \to$ uncertainty: $+0.0124$* to $+0.0309$*** across all eight
        specifications --- \emph{same sign} as US EPU.
        (Note: prior memory entry flagging a ``wrong-sign'' finding is
        stale; the rebuilt results align with US EPU in sign and
        magnitude. Memory to be updated.)
  \item \textbf{H25 (Table~\ref{tab:h25}).} $\log(\text{GPR})_t \to$
        uncertainty: mostly null, with a single $+0.0147$* in the
        CEO presentation column. Geopolitical risk does not
        translate cleanly into firm-level managerial language in our
        sample period.
\end{itemize}
References: Tables~\ref{tab:h11}, \ref{tab:h11_lag}, \ref{tab:h23},
\ref{tab:h24}, \ref{tab:h24b}, \ref{tab:h25}.\\
\textbf{Open:} DECISIONS.md \S3.5 --- whether to split this block into
(a) macro/political (H11, H11-Lag, H24, H24b, H25) and
(b) product-market (H23) as two separate sub-subsections.
\end{placeholder}

\subsubsection{Information-Environment Signals}
\label{sec:appII_info}
\begin{placeholder}
\textbf{Subnarrative.} Uncertain managerial speech correlates with
higher analyst disagreement and higher SEC scrutiny. These are
concurrent signals that the call transmits information to downstream
market participants.
\begin{itemize}
  \item \textbf{H5 (Table~\ref{tab:h5}).} Wang (2020) analyst forecast
        dispersion \texttt{DISP}: \texttt{UncAnsMgr} $+0.0002$*** ---
        statistically robust but economically very small.
  \item \textbf{H18 (Table~\ref{tab:h18}).} SEC comment letters (linear
        probability model): \texttt{UncPreMgr} $+0.0017$* to
        $+0.0020$*.
  \item \textbf{H18b (Table~\ref{tab:h18b}).} SEC comment letters
        (logit robustness): \texttt{UncPreMgr} $+0.0012$* to
        $+0.0018$**.
  \item \textbf{H21 (Table~\ref{tab:h21}).} SEC comment letter
        \emph{count}: \texttt{UncPreCEO} $+0.0134$** to $+0.0153$**.
\end{itemize}
References: Tables~\ref{tab:h5}, \ref{tab:h18}, \ref{tab:h18b},
\ref{tab:h21}.
\end{placeholder}

\subsubsection{Market Microstructure (Null)}
\label{sec:appII_micro}
\begin{placeholder}
\textbf{Subnarrative.} Managerial speech uncertainty predicts
corporate \emph{decisions} but does \emph{not} predict intraday
microstructure. Ten tables of Amihud illiquidity and bid-ask spread
tests over 3-day and 25-day post-call windows all return null. This
is an informative boundary: the effects we document run through real
corporate policy, not through trading-friction channels.
\begin{itemize}
  \item \textbf{H7 family (Tables~\ref{tab:h7}--\ref{tab:h7e}).}
        Amihud illiquidity: (a) 3-day change $[+1,+3]-[-3,-1]$;
        (b) 3-day level $[+1,+3]$; (c) BGT (2018) 25-day level
        $[0,+25]$; (d) BGT 25-day delta $[+1,+25]-[-25,-1]$;
        (e) BGT 51-day symmetric average $[-25,+25]$. All null on
        managerial IVs.
  \item \textbf{H14 family (Tables~\ref{tab:h14}--\ref{tab:h14e}).}
        Lee (2016) bid-ask spread: same five window variants as
        H7. All null.
\end{itemize}
References: Tables~\ref{tab:h7}--\ref{tab:h7e},
\ref{tab:h14}--\ref{tab:h14e}.\\
\textbf{Open:} DECISIONS.md \S3.6 --- whether to display all 10
tables separately, consolidate into one summary, or use two
family-level summaries.
\end{placeholder}

\clearpage

% =======================================================================
% Section V -- Conclusion
% =======================================================================
\section{Conclusion}
\label{sec:conclusion}
\begin{placeholder}
\textbf{Recap.}
\begin{enumerate}
  \item \textbf{Contribution.} (a) Decomposition of managerial speech
        uncertainty into answer-side (Q\&A, \texttt{UncAnsMgr}) and
        presentation-side (scripted, \texttt{UncPreMgr}) components;
        (b) extension from CEO-only to the full management team;
        (c) identification of multiple financial-conservatism channels
        in a single firm-quarter panel.
  \item \textbf{Main findings.} Precautionary cash rises; leverage
        drifts down at the one-quarter lead; dividends fall;
        repurchases fall; capex rises under competitive pressure as
        defensive preemption; R\&D is unresponsive.
  \item \textbf{Validation.} The measure responds in expected
        directions to US EPU, Global EPU, political risk, and
        product-market competition.
  \item \textbf{Boundary.} No intraday liquidity response --- the
        effect is corporate-policy, not microstructural.
  \item \textbf{Limitations.} Sample period 2002--2018 excludes the
        COVID shock and its aftermath. Current \texttt{UncAnsMgr}
        construction in the pipeline is broken pending the
        CFO / ManagerV2 rebuild (memory
        \texttt{project\_cfo\_mgrv2\_plan\_v2}); the present results
        should be replicated after that rebuild is executed.
  \item \textbf{Future work.} Channel moderation suites F2, F3, F5
        not yet built; cross-country extensions; alternative
        definitions of the manager category; interaction with
        institutional ownership and analyst coverage.
\end{enumerate}
\end{placeholder}

\clearpage

% =======================================================================
% References
% =======================================================================
\bibliographystyle{abbrvnat}
\bibliography{references}

\begin{placeholder}
\textbf{Populate \texttt{references.bib}.} Minimum load-bearing set:
\begin{itemize}
  \item Speech / disclosure: Davis, Ge, Matsumoto, \& Zhang (2015);
        Loughran \& McDonald (2011) uncertainty dictionary;
        Bushee, Gow, \& Taylor (2018, JAR); Larcker \& Zakolyukina
        (2012); \citet{DWZ2021}.
  \item Corporate conservatism: Bates, Kahle, \& Stulz (2009);
        Opler, Pinkowitz, Stulz, \& Williamson (1999); Myers (1984);
        Frank \& Goyal (2003); Lintner (1956); Brav, Graham, Harvey,
        \& Michaely (2005); Hoberg \& Prabhala (2009); Jiang (2021).
  \item Real options and investment under uncertainty: Dixit \&
        Pindyck (1994); Bloom (2014); Grenadier (2002); Aguerrevere
        (2009).
  \item Financing-choice margins: Chang, Dasgupta, \& Hilary (2006);
        Hoberg \& Maksimovic.
  \item Uncertainty IVs and validation: Baker, Bloom, \& Davis (EPU);
        Davis (GEPU); Caldara \& Iacoviello (GPR); Hassan, Hollander,
        van Lent, \& Tahoun (PRisk); Hoberg \& Phillips (TSIMM).
  \item Microstructure: Amihud (2002); Lee (2016).
  \item Dispersion / expectations: Wang (2020).
\end{itemize}
\end{placeholder}

\clearpage

% =======================================================================
% Appendix: Variable Definition and Measurement
% =======================================================================
\appendix
\section{Variable Definition and Measurement}
\label{app:vardef}

\subsection{Dependent Variables}
\label{app:vardef_dv}
\begin{placeholder}
Populate \emph{verbatim} from \texttt{findings.txt} DV formula block.
Each entry must spell out every raw database item with its source
(memory rule \texttt{feedback\_dv\_formula\_documentation}). Required
entries, in section order:
\begin{itemize}
  \item \texttt{CashRatio} (H1) --- Compustat \texttt{cheq}/\texttt{atq}.
  \item \texttt{Leverage} (H4a) --- (\texttt{dlcq} + \texttt{dlttq}) /
        \texttt{atq}.
  \item \texttt{DebtToCapital} (H4b) --- (\texttt{dlcq} +
        \texttt{dlttq}) / (\texttt{dlcq} + \texttt{dlttq} +
        \texttt{ceqq}).
  \item \texttt{PayoutRatio} (H12) --- de-cumulated
        \texttt{dvy} / \texttt{ibq}.
  \item \texttt{DivPayer} (H12b) --- indicator of positive quarterly
        dividends (Hoberg-Prabhala 2009 analog).
  \item \texttt{Capex} (H13, H13.1, H13.2) --- de-cumulated
        \texttt{capxy} / \texttt{atq}.
  \item \texttt{RDSales} (H16) --- \texttt{xrdy} / \texttt{saley}
        (Jiang 2021).
  \item \texttt{RepurchaseIntensity} (H17) --- de-cumulated
        \texttt{prstkcy} / \texttt{atq}.
\end{itemize}
\end{placeholder}

\subsection{Independent Variables}
\label{app:vardef_iv}
\begin{placeholder}
\begin{itemize}
  \item \texttt{UncAnsMgr} --- manager answer-side (Q\&A) uncertainty,
        DWZ (2021) persistent-style construction with Loughran-McDonald
        uncertainty word list. Built via \texttt{linguistic\_engine}
        over the Q\&A segment for non-CEO executives.
  \item \texttt{UncPreMgr} --- manager presentation-side (scripted)
        uncertainty, same construction over the opening presentation
        segment.
  \item \texttt{UncAnsCEO}, \texttt{UncPreCEO} --- CEO analogs, using
        Execucomp \texttt{ceoann} annual flag.
  \item \textbf{Moderators:}
        \texttt{BelowIG}, \texttt{Unrated} (H1.2);
        \texttt{z\_log\_TotalSimilarity}, \texttt{HighTSIMM}
        (H1.1, H13.1).
\end{itemize}
\end{placeholder}

\subsection{Controls}
\label{app:vardef_controls}
\begin{placeholder}
\textbf{Base:} \texttt{lnAssets}, \texttt{TobinsQ}, \texttt{ROA},
\texttt{Leverage}, \texttt{CashRatio}, \texttt{Capex},
\texttt{DivDummy}, \texttt{sCFO}, \texttt{Lagged\_DV}.\\
\textbf{Extended:} \texttt{SalesGrowth}, \texttt{RDSales},
\texttt{CashFlowAt}, \texttt{DailyVola}.\\
\textbf{Bad-control exclusions.} Per DV, controls that share the
DV's numerator concept are dropped (e.g., \texttt{CashRatio} and
\texttt{CashFlowAt} not used as controls when \texttt{CashRatio}
is the DV). The full exclusion map is enforced by the shared
runner configuration in \texttt{src/f1d/econometric/}.
\end{placeholder}

\clearpage

% =======================================================================
% Tables
% =======================================================================
\section{Tables}
\label{app:tables}

\begin{placeholder}
\textbf{Table-fragment strategy.} The authoritative
\texttt{all\_tables.tex} is a 2{,}350-line monolith generated by
\texttt{generate\_all\_tables.py}. To use \texttt{\textbackslash input\{\}}
from this draft, we will need a future sub-task to split that
monolith into per-suite fragments, e.g., \texttt{tables/tab\_h1.tex},
\texttt{tables/tab\_h1\_2.tex}, \ldots. Until that split is done,
the \texttt{\textbackslash input\{\}} hooks below are left commented
out. Table-number references (\texttt{\textbackslash ref\{tab:h1\}}
etc.) will resolve once the fragments are imported.
\end{placeholder}

% -----------------------------------------------------------------------
% Main body tables
% -----------------------------------------------------------------------
\subsection*{Table 1. Summary Statistics}
\begin{placeholder}
Generate from \texttt{findings.txt} or a dedicated summary script.
Main sample: 112{,}968 calls, 2{,}429 firms, 2002 Q1--2018 Q4.
Report mean / SD / N for all DVs, all four speech-uncertainty IVs,
and the full control set. Break out by main sample (ex fin\&util).
\end{placeholder}

\subsection*{Table 2. H1 --- Speech Uncertainty and Cash Holdings}
% \input{tables/tab_h1.tex}
\begin{placeholder}
H1 placeholder --- \texttt{tables/tab\_h1.tex}.
\end{placeholder}

\subsection*{Table 3. H1.2 --- Rating-Channel Moderation of Cash Holdings}
% \input{tables/tab_h1_2.tex}
\begin{placeholder}
H1.2 placeholder --- three-category financial constraint.
\end{placeholder}

\subsection*{Table 4. H4a --- Book Leverage}
% \input{tables/tab_h4a.tex}
\begin{placeholder}
H4a placeholder.
\end{placeholder}

\subsection*{Table 5. H4b --- Debt-to-Capital}
% \input{tables/tab_h4b.tex}
\begin{placeholder}
H4b placeholder.
\end{placeholder}

\subsection*{Table 6. H12 --- Quarterly Payout Ratio}
% \input{tables/tab_h12.tex}
\begin{placeholder}
H12 placeholder.
\end{placeholder}

\subsection*{Table 7. H12b --- Dividend Payer Indicator}
% \input{tables/tab_h12b.tex}
\begin{placeholder}
H12b placeholder --- Hoberg-Prabhala 2009 analog.
\end{placeholder}

\subsection*{Table 8. H13 --- Capital Expenditure}
% \input{tables/tab_h13.tex}
\begin{placeholder}
H13 placeholder.
\end{placeholder}

\subsection*{Table 9. H13.1 --- Product-Similarity Moderation of Capex}
% \input{tables/tab_h13_1.tex}
\begin{placeholder}
H13.1 placeholder --- TSIMM moderation interaction.
\end{placeholder}

\subsection*{Table 10. H16 --- R\&D Investment Intensity}
% \input{tables/tab_h16.tex}
\begin{placeholder}
H16 placeholder --- Jiang 2021 RDSales.
\end{placeholder}

\subsection*{Table 11. H17 --- Repurchase Intensity}
% \input{tables/tab_h17.tex}
\begin{placeholder}
H17 placeholder.
\end{placeholder}

% -----------------------------------------------------------------------
% Appendix I tables (main-test robustness)
% -----------------------------------------------------------------------
\subsection*{Appendix I Tables (Main-Test Robustness)}
\begin{placeholder}
A.I.1: H1.1 product-similarity moderation of cash (null interaction).\\
A.I.2: H1.1b binary TSIMM variant of H1.1 (null interaction).\\
A.I.3: H13.2 capex lead horizons $t+1$ through $t+4$.\\
A.I.4: H19b external vs.\ internal financing (CDH 2006).\\
A.I.5: H20b debt vs.\ equity choice (CDH 2006).
\end{placeholder}

% -----------------------------------------------------------------------
% Appendix II tables (validity, drivers, boundary tests)
% -----------------------------------------------------------------------
\subsection*{Appendix II Tables (Validity, Drivers, and Boundary Tests)}
\begin{placeholder}
\textbf{A.II.1 External drivers of managerial speech uncertainty:}
H11 (PRisk contemporaneous), H11-Lag (lagged PRisk), H23 (TSIMM
$\to$ uncertainty, reverse), H24 (US EPU), H24b (GEPU), H25 (GPR).\\
\textbf{A.II.2 Information environment signals:}
H5 (Wang 2020 DISP), H18 (SEC comment letter LPM), H18b (SEC comment
letter logit), H21 (SEC comment letter count).\\
\textbf{A.II.3 Market microstructure (null):}
H7 family --- 3-day $\Delta$Amihud, 3-day Amihud level, BGT 25-day
level, BGT 25-day $\Delta$, BGT 51-day symmetric average.
H14 family --- 3-day $\Delta$DSPREAD, 3-day DSPREAD level, BGT 25-day
variants. Ten tables total; display strategy pending DECISIONS.md
\S3.6.
\end{placeholder}

\end{document}
